\chapter*{Cómo usar este manual}
\addcontentsline{toc}{chapter}{Cómo usar este manual}

Este manual explica, paso a paso, cómo se construyó \textbf{FuegoNet}: una red
neuronal de detección de objetos con la misma arquitectura que YOLO26, escrita
desde cero en PyTorch, para reconocer \textbf{fuego} y \textbf{humo} en las
imágenes de la cámara del Tarot 650 durante la búsqueda de incendios forestales.

Hasta ahora la simulación usaba YOLO26n de Ultralytics: se descargan pesos
entrenados en COCO (80 clases de objetos cotidianos) y se ajustan a nuestras dos
clases. Funciona muy bien, pero la red es una caja negra: no se ve qué hace por
dentro, no se puede modificar con libertad y, en un trabajo de investigación, es difícil
justificar cada decisión. Aquí abrimos esa caja. Cada pieza de YOLO26 se
programa a mano, se prueba por separado y se explica con sus ecuaciones. Al final
se entrena la red completa \emph{sin pesos preentrenados} y se compara con
YOLO26n usando el mismo código de medición.

\section*{Qué necesitas saber antes}
\begin{itemize}
  \item Python básico (funciones, clases, listas) y saber correr un comando en la terminal.
  \item La idea de neurona, convolución, IoU y mAP. Si no la tienes, lee primero
        el capítulo «Redes neuronales y YOLO sin complicaciones» del libro de la
        simulación (\cod{doc/manual}). Aquí se repasa lo indispensable.
  \item No hace falta saber PyTorch: cada vez que aparece algo nuevo se explica.
\end{itemize}

\section*{Cómo está organizado}
\begin{description}
  \item[Parte I] Qué es YOLO26, qué tiene de nuevo y cómo preparar el entorno y los datos.
  \item[Parte II] Los bloques de la red (convoluciones, C3k2, SPPF, atención) y cómo se
        ensamblan en columna vertebral, cuello y cabeza.
  \item[Parte III] Cómo aprende: la asignación de etiquetas (TAL y STAL), la función
        de pérdida de dos cabezas, los aumentos de datos y el ciclo de entrenamiento
        con el optimizador MuSGD.
  \item[Parte IV] Cómo se mide, cómo se compara con YOLO26n y cómo se conecta a ROS 2
        para que el dron la use en vuelo.
\end{description}

Cada capítulo avanza en \textbf{pasos numerados}. Al final de cada paso hay un
recuadro «Lo que debes ver» con la salida esperada; si lo que ves es distinto,
no sigas: algo quedó mal en ese paso.

\section*{Convenciones}
\begin{itemize}
  \item Los comandos se escriben así y se corren desde \cod{~/ros2_ws/src/cuadrirotor_sim}:
\begin{lstlisting}[style=bash]
~/.venvs/yolo/bin/python -m cuadrirotor_sim.red_fuego.modelo
\end{lstlisting}
  \item Las rutas, archivos y nombres de variables van en \cod{monoespaciado}.
  \item El código completo está en \cod{cuadrirotor_sim/red_fuego/} y en el
        apéndice~\ref{ap:codigo}. En el texto solo aparecen las partes importantes.
  \item Los nombres del código están en español (\cod{Cuello}, \cod{Atencion},
        \cod{AsignadorTAL}) para que se lean junto con el texto. Entre paréntesis se
        da el nombre en inglés que vas a encontrar en los artículos.
\end{itemize}

\begin{nota}
Todos los números de este manual (parámetros, tiempos, mAP) salieron de correr el
código en la computadora del laboratorio: Ubuntu, PyTorch 2.14 con CUDA 13 y una
tarjeta NVIDIA RTX~5060~Ti de 16~GB. En otra máquina los tiempos cambian, pero
las métricas deben salir parecidas.
\end{nota}
